\documentclass[11pt, a4paper]{article}

\usepackage[T1]{fontenc}
\usepackage[utf8]{inputenc}
\usepackage{tgpagella}
\usepackage{tgcursor}
\usepackage{microtype}
\usepackage[spanish, es-noshorthands]{babel}
\addto\captionsspanish{\renewcommand{\tablename}{Tabla}}

\usepackage{amsmath, amssymb, amsfonts, bm}
\usepackage[table, xcdraw]{xcolor}
\usepackage{graphicx}
\usepackage{booktabs}
\usepackage{tabularx}
\usepackage[most]{tcolorbox}
\usepackage[margin=2.5cm]{geometry}
\usepackage{fancyhdr}

\usepackage{listings}
\lstset{
    language=Python,
    basicstyle=\ttfamily\small,
    backgroundcolor=\color{gray!5},
    frame=single,
    rulecolor=\color{gray!40},
    keywordstyle=\color{blue}\bfseries,
    stringstyle=\color{red!70!black},
    commentstyle=\color{green!60!black}
}

\pagestyle{fancy}
\fancyhf{}
\setlength{\headheight}{16pt}
\lhead{\textbf{UCB Sede Tarija} | Ing. Mecatrónica}
\chead{Robótica (IMT-342)}
\rhead{Solucionario Docente}
\lfoot{Prof: M.Sc. Ing. Bernardo Quiroga Turdera}
\rfoot{Página \thepage}
\renewcommand{\headrulewidth}{0.4pt}
\renewcommand{\footrulewidth}{0.4pt}

\definecolor{ucbblue}{RGB}{0, 51, 102}

\begin{document}

\begin{center}
    \Large\textbf{\textcolor{ucbblue}{UNIVERSIDAD CATÓLICA BOLIVIANA "SAN PABLO"}}\\
    \vspace{0.2cm}
    \normalsize \textbf{SOLUCIONARIO DEL INSTRUCTOR (No Distribuir)} \\
    Asignatura: Robótica Industrial (IMT-342) \\
    Sesión: W05-S01 - Cinemática Directa
\end{center}

\vspace{0.4cm}
\begin{tcolorbox}[colback=red!5!white, colframe=red!75!black, title=\textbf{Notas Diagnósticas para el Docente}]
    Esta hoja contiene las respuestas esperadas y heurísticas de corrección basadas en las dificultades anticipadas de la sesión[cite: 1, 2].
\end{tcolorbox}
\vspace{0.4cm}

\section*{Solución Ejercicio 1[cite: 1]}
1. \textbf{Matriz Genérica:} Sustituyendo en la fórmula D-H estándar:
\begin{equation*}
    {}^{i-1}T_i = \begin{bmatrix}
    \cos q_i & -\sin q_i & 0 & 0.25\cos q_i \\
    \sin q_i & \cos q_i & 0 & 0.25\sin q_i \\
    0 & 0 & 1 & 0 \\
    0 & 0 & 0 & 1
    \end{bmatrix}
\end{equation*}
2. \textbf{Evaluación numérica ($90^\circ$):}
\begin{equation*}
    {}^{i-1}T_i = \begin{bmatrix}
    0 & -1 & 0 & 0 \\
    1 & 0 & 0 & 0.25 \\
    0 & 0 & 1 & 0 \\
    0 & 0 & 0 & 1
    \end{bmatrix}
\end{equation*}
3. \textbf{Interpretación:} La última columna indica un desplazamiento de $0.25\text{ m}$ exclusivamente a lo largo del eje $y_{i-1}$ debido al giro de $90^\circ$. \\
\textit{Diagnóstico:} Si el estudiante deja $[0.25, 0, 0]^T$ en la última columna, significa que no aplicó el operador de rotación al vector de desplazamiento del eslabón.

\section*{Solución Ejercicio 2[cite: 1]}
1. \textbf{Matrices:} Para $q_1=0^\circ, q_2=90^\circ, q_3=-90^\circ$:
\begin{equation*}
    {}^{0}T_1 = \begin{bmatrix} 1&0&0&0.30 \\ 0&1&0&0 \\ 0&0&1&0 \\ 0&0&0&1 \end{bmatrix}, \quad
    {}^{1}T_2 = \begin{bmatrix} 0&-1&0&0 \\ 1&0&0&0.25 \\ 0&0&1&0 \\ 0&0&0&1 \end{bmatrix}, \quad
    {}^{2}T_3 = \begin{bmatrix} 0&1&0&0 \\ -1&0&0&-0.15 \\ 0&0&1&0 \\ 0&0&0&1 \end{bmatrix}
\end{equation*}
2. \textbf{Secuencia:} ${}^{0}T_3 = {}^{0}T_1 {}^{1}T_2 {}^{2}T_3$.
3. \textbf{Resultado Numérico:} 
\begin{equation*}
    {}^{0}T_3 = \begin{bmatrix} 1&0&0&0.45 \\ 0&1&0&0.25 \\ 0&0&1&0 \\ 0&0&0&1 \end{bmatrix}
\end{equation*}
4. \textbf{Posición:} $p_x = 0.45\text{ m}, p_y = 0.25\text{ m}, p_z = 0\text{ m}$.
5. \textbf{Verificación:} La magnitud máxima planar teórica es $0.30+0.25+0.15 = 0.70\text{ m}$. El módulo actual es $\sqrt{0.45^2 + 0.25^2} \approx 0.515 < 0.70\text{ m}$. Es una posición geométricamente alcanzable.

\section*{Solución Ejercicio 3 (ABB IRB 120)[cite: 1]}
1. \textbf{Offset explícito:} La matriz ${}^{1}T_2$ presenta senos y cosenos evaluados en $-\pi/2$, reflejando $\theta_2 = 0 - \pi/2 = -\pi/2$. (Nótese que $(1,1)$ es 0 y $(2,1)$ es -1).
2. \textbf{Cadena:} ${}^{0}T_6 = {}^{0}T_1 {}^{1}T_2 {}^{2}T_3 {}^{3}T_4 {}^{4}T_5 {}^{5}T_6$.
3. \textbf{Resultado Final:}
\begin{equation*}
    {}^{0}T_6 = \begin{bmatrix}
    0 & 0 & 1 & 0.374 \\
    0 & -1 & 0 & 0 \\
    1 & 0 & 0 & 0.630 \\
    0 & 0 & 0 & 1
    \end{bmatrix}
\end{equation*}
4. \textbf{Extracción:} $\mathbf{p} = [0.374, 0, 0.630]^T\text{ m}$. Ejes: $x_6 \parallel +z_0$, $y_6 \parallel -y_0$, $z_6 \parallel +x_0$.
5. \textbf{Explicación del Cero:} Las matrices provienen de una convención D-H arbitraria (modelo analítico). El "cero" de fábrica depende de los topes mecánicos y la calibración del fabricante, los cuales pueden tener desfases angulares constantes adicionales[cite: 1].

\section*{Script de Verificación Computacional (Python)[cite: 1]}
\begin{lstlisting}
import numpy as np

def dh_std(theta, d, a, alpha):
    ct = np.cos(theta); st = np.sin(theta)
    ca = np.cos(alpha); sa = np.sin(alpha)
    return np.array([
        [ct, -st*ca,  st*sa, a*ct],
        [st,  ct*ca, -ct*sa, a*st],
        [0,      sa,     ca,    d],
        [0,       0,      0,    1]
    ], dtype=float)

T01 = dh_std(0, 0.290, 0.0, -np.pi/2)
T12 = dh_std(-np.pi/2, 0.0, 0.270, 0.0)
T23 = dh_std(0, 0.0, 0.070, -np.pi/2)
T34 = dh_std(0, 0.302, 0.0, np.pi/2)
T45 = dh_std(0, 0.0, 0.0, -np.pi/2)
T56 = dh_std(0, 0.072, 0.0, 0.0)

T06 = T01 @ T12 @ T23 @ T34 @ T45 @ T56
print(np.round(T06, 3))
\end{lstlisting}

\end{document}
